% Introduction: Presenting the phenomenon without attempting to fully explain it
% Target: Nature or Nature Materials
% Revised after Round 5 review.

The link between local structure and ion transport in solid-state electrolytes (SSEs) remains poorly understood. SSEs are a promising class of materials for next-generation battery technologies, offering improved safety and energy density compared to liquid electrolytes \cite{ref1}. Among them, sodium superionic conductors based on closo-borate anions (e.g., Na$_2$B$_{10}$H$_{10}$ and mixed-anion systems) exhibit exceptionally high ionic conductivities at moderate temperatures \cite{ref2,ref3}, yet the microscopic mechanisms—and the role of grain-scale dynamics—are still unclear.

X-ray photon correlation spectroscopy (XPCS) provides a unique probe of slow dynamics in condensed matter by measuring temporal correlations in coherent X-ray scattering \cite{ref4,ref5}. In Bragg geometry, XPCS is sensitive to phase-coherent fluctuations in strain, displacement fields, and domain structure, making it well-suited to investigate collective dynamics in crystalline materials \cite{ref6}. The two-time correlation function $\ttc(t_1, t_2)$—correlations between speckle patterns at different experimental times—reveals both intrinsic temporal structure and non-stationary evolution, offering more than conventional autocorrelation analysis \cite{ref7,ref8}.

Here, we apply XPCS to grain-scale dynamics in closo-borate SSEs, comparing the prototypical Na$_2$B$_{10}$H$_{10}$ (B10) system with mixed-anion compositions. B10 exhibits a comparatively simpler correlation structure, whereas in mixed-anion systems we observe pronounced checkerboard-like patterns in $\ttc(t_1, t_2)$—spatially structured fluctuations that vary strongly with reciprocal-space position ($\qvec$-vector), with the most pronounced structure in certain regions of the Bragg peak. We do not assign a unique structural or mechanistic origin to these patterns in this work.

A particularly striking observation is the appearance of oscillatory correlations along anti-diagonal cuts through the two-time correlation function. These oscillations occur at fixed mean experimental time $T = (t_1 + t_2)/2$ as a function of lag time $\tauval = t_2 - t_1$, with characteristic periods on the order of tens of seconds at the time resolution used here. They cannot be explained by aging, instrumental drift, or simple diffusive relaxation: aging would appear along the diagonal, drift breaks symmetry in $t_1$--$t_2$, and diffusion yields monotonic decay in $\tauval$. We have verified that these features are not introduced by the data processing or detector response. The oscillations are consistent with the system revisiting similar microscopic configurations in a phase-coherent manner. Their frequency and amplitude vary systematically with reciprocal-space position, suggesting spatially heterogeneous dynamics; similar qualitative behaviour is observed across multiple samples and runs, as shown below.

Temperature-dependent measurements from $\sim$80~K to room temperature reveal further complexity. The two-time maps exhibit a marked change in structure near $\sim$120~K and at room temperature, indicating distinct dynamical regimes. Quantitative analysis yields characteristic frequencies or timescales associated with these dynamics.

The physical origin of these dynamics is not yet fully established. We focus here on demonstrating that XPCS reveals collective, grain-scale processes in closo-borate SSEs that are not accessible through conventional structural or transport measurements. This work establishes XPCS as a powerful tool for investigating dynamics in solid-state battery materials and opens new avenues for understanding the relationship between local structure, local ordering, and ionic conductivity. Mechanistic interpretation and modelling are left for future work.
